
\documentclass[11pt,a4paper]{article}

% Packages
\usepackage{graphicx}
\usepackage[utf8]{inputenc}
\usepackage[T1]{fontenc}
\usepackage{geometry}
\usepackage{booktabs}
\usepackage{array}
\usepackage{float}
\usepackage{amsmath}
\usepackage{enumitem}
\usepackage{microtype}
\usepackage{hyperref}
\usepackage{xcolor}

\geometry{margin=1in}
\setlength{\parindent}{0pt}
\setlength{\parskip}{7pt}
\renewcommand{\arraystretch}{1.2}

\title{\textbf{Question 1: Archetype-Level Descriptive Statistics}\\
\large Results and Interpretations}
\author{}
\date{}

\begin{document}
\maketitle

\section{Research Question}

Generate archetype-level descriptive statistics for
age, height, weight, and BMI across the various
functional roles of Indian Olympic athletes.

The analysis examines differences in central
tendency, variability, distributional characteristics,
and sex-wise anthropometric profiles across seven
functional archetypes.

\section{Sample Composition and Archetype Distribution}

The dataset contains 737 athlete-event-edition
observations representing 522 unique athletes,
classified into seven functional archetypes.

\begin{table}[H]
\centering
\caption{Distribution of Observations Across Archetypes}
\begin{tabular}{lrr}
\toprule
\textbf{Archetype} &
\textbf{Observations} &
\textbf{Unique Athletes} \\
\midrule
Combat             & 110 & 84  \\
Endurance          & 56  & 44  \\
Precision          & 97  & 61  \\
Racket Sports      & 61  & 36  \\
Speed \& Power     & 139 & 104 \\
Team Sports        & 245 & 173 \\
Technical / Mixed  & 29  & 24  \\
\midrule
\textbf{Total}     & \textbf{737} & \textbf{522} \\
\bottomrule
\end{tabular}
\end{table}

\textit{Note:} Unique-athlete counts within
archetypes are not additive because some athletes
may belong to multiple archetypes.

\subsection*{Interpretation}

Team Sports accounts for the largest share of
observations (33.2\%), followed by Speed \& Power
(18.9\%) and Combat (14.9\%).

Technical / Mixed has the smallest representation,
contributing only 3.9\% of observations.

This unequal distribution implies that descriptive
estimates for larger archetypes are based on more
observations, whereas estimates for smaller groups
may be more sensitive to individual athletes.

Additionally, some athletes appear in multiple
Olympic events or editions. Therefore, the results
describe recorded observations rather than an
equally weighted sample of independent athletes.


\section{Comparison of Mean Anthropometric Characteristics}

\begin{table}[H]
\centering
\small
\caption{Mean Anthropometric Characteristics by Archetype}
\resizebox{\textwidth}{!}{%
\begin{tabular}{lrrrr}
\toprule
\textbf{Archetype} &
\textbf{Age (years)} &
\textbf{Height (cm)} &
\textbf{Weight (kg)} &
\textbf{BMI} \\
\midrule
Combat            & 24.10 & 169.64 & 68.85 & 23.69 \\
Endurance         & 26.21 & 172.54 & 61.04 & 20.36 \\
Precision         & 29.21 & 171.80 & 71.57 & 24.07 \\
Racket Sports     & 26.20 & 176.66 & 70.98 & 22.66 \\
Speed \& Power    & 24.90 & 171.25 & 68.06 & 22.85 \\
Team Sports       & 24.74 & 171.80 & 66.20 & 22.41 \\
Technical / Mixed & 28.90 & 172.59 & 65.76 & 22.03 \\
\bottomrule
\end{tabular}%
}
\end{table}

\subsection*{Interpretation}

Precision athletes have the highest average age,
weight, and BMI, while Racket Sports athletes are
the tallest on average.

Endurance athletes exhibit the lowest mean weight
and BMI, despite having an average height
comparable to most other archetypes.

These observations suggest that physical profiles
differ across functional categories.

However, these are descriptive comparisons and
cannot establish that the demands of a particular
sporting category cause the observed physical
characteristics.


\section{Age Distribution Across Functional Archetypes}

\subsection*{Interpretation}

Precision athletes have the highest mean age
(29.21 years), closely followed by Technical /
Mixed athletes (28.90 years).

Combat athletes have the lowest mean age
(24.10 years), followed by Team Sports
(24.74 years) and Speed \& Power (24.90 years).

The difference between the highest and lowest
archetype means is approximately
\textbf{5.11 years}, indicating noticeable
variation in age profiles across the groups.

Precision also exhibits the greatest age variability,
with a standard deviation of 8.57 years and an
observed age range of 16--66 years.

This indicates a relatively heterogeneous age
distribution, including considerably older
participants.

By contrast, Team Sports has the smallest age
standard deviation (3.60 years), suggesting
that its observations are concentrated within
a narrower age range.

The median age for Precision is 28 years,
below its mean of 29.21 years.

Together with its high maximum age, this
suggests a right-skewed distribution influenced
by older athletes.

\textbf{Key Finding:}
Precision and Technical / Mixed athletes tend
to be older, while Combat and Team Sports
exhibit younger and relatively concentrated
age profiles.


\section{Height Distribution Across Functional Archetypes}

\subsection*{Interpretation}

Racket Sports athletes have the highest average
height (176.66 cm), approximately 7.02 cm greater
than Combat athletes, who have the lowest mean
height (169.64 cm).

Endurance, Precision, Speed \& Power, Team Sports,
and Technical / Mixed have mean heights between
approximately 171 and 173 cm.

This indicates relatively small differences in
average stature among these five groups.

Speed \& Power has the largest height standard
deviation (11.54 cm), suggesting substantial
variation in body stature within this archetype.

This is consistent with the possibility that
its constituent events include athletes with
different physical requirements.

Technical / Mixed has the smallest height
standard deviation (7.80 cm), although its
relatively small sample size should be considered
when comparing variability.

The median height of Racket Sports athletes
(178 cm) also exceeds that of the other groups.

This indicates that its relatively high mean
is not merely the result of a few exceptionally
tall observations.

\textbf{Key Finding:}
Racket Sports shows a comparatively taller
physical profile, whereas height differences
among most other archetypes are relatively modest.


\section{Weight Distribution Across Functional Archetypes}

\subsection*{Interpretation}

Precision athletes have the highest mean weight
(71.57 kg), closely followed by Racket Sports
(70.98 kg).

Endurance athletes have the lowest mean weight
(61.04 kg).

The difference between the highest and lowest
mean weights is approximately
\textbf{10.53 kg}, demonstrating considerable
variation in body mass across functional
archetypes.

Speed \& Power exhibits the greatest weight
variability, with a standard deviation of
19.15 kg and weights ranging from 46 to 125 kg.

Its mean weight (68.06 kg) is substantially
greater than its median (61 kg), indicating
a right-skewed distribution in which heavier
observations increase the mean.

Combat athletes also show substantial weight
variability (SD = 16.47 kg), which may reflect
differences in body size and weight-category
participation.

By contrast, Team Sports has a comparatively
low weight standard deviation (7.89 kg),
suggesting a more concentrated weight
distribution.

Endurance athletes maintain the lowest average
weight despite having an average height
of 172.54 cm, similar to several other
archetypes.

\textbf{Key Finding:}
Endurance athletes display a relatively lighter
physical profile, whereas Speed \& Power and
Combat exhibit greater diversity in body mass.


\section{BMI Distribution Across Functional Archetypes}

\subsection*{Interpretation}

Precision athletes have the highest mean BMI
(24.07), followed by Combat (23.69).

Endurance athletes have the lowest mean
BMI (20.36).

The difference between the highest and lowest
mean BMI values is approximately
\textbf{3.71 kg/m$^2$}.

This indicates descriptive variation in body
weight relative to height.

Endurance has the lowest BMI standard deviation
(1.82), suggesting that its observations are
clustered relatively closely around a lower BMI.

Racket Sports and Team Sports also exhibit
comparatively concentrated BMI distributions.

Precision has the largest BMI standard deviation
(4.32), followed by Combat (4.02) and
Speed \& Power (3.91).

This indicates greater variation in
height-adjusted body mass within these groups.

The maximum recorded BMI in Combat is
approximately 39.18, considerably above its
group median of 22.94.

Similarly, Speed \& Power and Precision contain
relatively high BMI observations.

These values contribute to the observed upper
tails and greater dispersion.

Importantly, BMI cannot distinguish muscle
mass from body fat.

Therefore, a higher BMI among athletes should
not be interpreted directly as greater adiposity
or poorer physical fitness.

\textbf{Key Finding:}
Endurance athletes show the lowest and most
consistent BMI profile, while Precision,
Combat, and Speed \& Power exhibit greater
heterogeneity in height-adjusted body mass.


\section{Sex-Wise Comparisons}

Sex-wise analysis is important because the
relative proportions of male and female
observations vary substantially between
archetypes.

\begin{table}[H]
\centering
\small
\caption{Sex-Wise Comparison of Mean Height and BMI}
\resizebox{\textwidth}{!}{%
\begin{tabular}{lrrrr}
\toprule
\textbf{Archetype} &
\textbf{Female Height} &
\textbf{Male Height} &
\textbf{Female BMI} &
\textbf{Male BMI} \\
\midrule
Combat
& 158.88 & 170.48 & 22.71 & 23.77 \\
Endurance
& 162.00 & 175.72 & 18.63 & 20.89 \\
Precision
& 162.86 & 175.43 & 21.90 & 24.95 \\
Racket Sports
& 167.95 & 181.56 & 22.06 & 23.00 \\
Speed \& Power
& 165.78 & 175.79 & 21.53 & 23.94 \\
Team Sports
& 159.94 & 172.62 & 21.22 & 22.49 \\
Technical / Mixed
& 167.50 & 174.52 & 21.02 & 22.42 \\
\bottomrule
\end{tabular}%
}
\end{table}

\subsection{Age Comparison}

Within Precision, male observations have a
substantially higher mean age (30.96 years)
than female observations (24.89 years).

A similar pattern appears in Technical / Mixed,
where males average 30.19 years compared with
25.50 years for females.

Conversely, Speed \& Power has a slightly higher
mean age among females (25.21 years) than males
(24.64 years).

Endurance has almost identical mean ages for
males and females.

These results indicate that differences in age
composition between sexes are not uniform
across archetypes.


\subsection{Height Comparison}

Male observations have a higher mean height
than female observations across all seven
archetypes.

The largest mean difference occurs in Endurance
(13.72 cm), closely followed by Racket Sports
(13.61 cm).

Racket Sports has the highest mean height for
both males (181.56 cm) and females (167.95 cm).

This indicates that its relatively tall profile
persists when the sexes are analyzed separately.


\subsection{Weight Comparison}

Male observations have a higher mean weight
across every archetype.

The largest sex-wise mean weight difference
occurs in Precision, where males average
77.03 kg compared with 58.11 kg among females.

For both sexes, Endurance has the lowest
mean weight: 64.70 kg for males and 48.92 kg
for females.

This shows that the relatively low weight
of Endurance athletes is not solely an
artifact of the overall male-to-female
composition.


\subsection{BMI Comparison}

Male observations have a higher mean BMI
in every archetype.

The largest difference appears in Precision,
where mean BMI is 24.95 for males and
21.90 for females.

Among female observations, Endurance has
the lowest mean BMI (18.63), while Combat
has the highest (22.71).

Among males, Endurance is again the lowest
(20.89), while Precision is the highest
(24.95).

Thus, the relatively low BMI of Endurance
athletes persists for both sexes.

This strengthens the descriptive observation
that this archetype has a lighter
height-adjusted body-mass profile.


\subsection*{Important Limitation}

Sex representation is highly unequal in
several groups.

For example, Team Sports has 229 male
observations but only 16 female observations,
and Combat has 102 male observations
compared with eight female observations.

These imbalances limit the reliability of
direct comparisons, particularly within
the smaller female subgroups.


\section{Distributional Patterns from Histograms and Boxplots}

The histogram and boxplot comparisons provide
additional information beyond differences
in group means.

\subsection{Age}

Precision has the widest age distribution,
including observations substantially older
than those in most other groups.

Team Sports and Combat show comparatively
concentrated age distributions.

The presence of older observations in Precision
contributes to its higher mean and standard
deviation.


\subsection{Height}

Most archetypes show considerable overlap in
their height distributions, even though their
means differ.

Racket Sports has a relatively high median
height, while Speed \& Power exhibits the
widest interquartile range (17 cm).

This indicates substantial variation within
the middle 50\% of its height observations.


\subsection{Weight}

Speed \& Power and Combat exhibit substantial
variation in weight.

In Speed \& Power, the difference between the
mean and median is particularly pronounced,
suggesting that heavier athletes contribute
to an extended upper tail.

Team Sports has a narrower interquartile
range (10 kg), indicating greater concentration
in the central portion of its weight distribution.


\subsection{BMI}

Endurance has a relatively small BMI
interquartile range (2.30), while Technical /
Mixed has an even smaller IQR (2.24).

Combat has the largest BMI interquartile
range (6.31), indicating substantial differences
in height-adjusted weight even among the
central 50\% of observations.

The presence of extreme BMI values in certain
archetypes reinforces the need to examine
distributions rather than relying exclusively
on mean values.

Overall, the archetypes differ not only in
their average physical characteristics but
also in the extent of within-group variation.


\section{Height--Weight Relationship}

The height--weight scatterplot examines whether
taller athletes generally tend to have higher
body weight.

The pooled observations show a moderately strong
positive Pearson correlation:

\[
\boxed{r = 0.698}
\]

This indicates that taller athletes generally
tend to be heavier, although substantial
variation remains.

The strength of this association varies
by archetype.

Endurance exhibits a particularly strong
positive correlation:

\[
r_{\text{Endurance}} = 0.882
\]

Precision shows a weaker, though still
positive, correlation:

\[
r_{\text{Precision}} = 0.635
\]

These differences suggest that the relationship
between height and weight is not identical
across functional groups.

However, these correlations do not establish
causal relationships.

Pooled associations may also reflect differences
in sex composition and sporting events.


\section{Age--BMI Relationship}

The overall Pearson correlation between
age and BMI is approximately:

\[
\boxed{r = 0.240}
\]

This indicates a weak positive linear
association.

However, the relationship varies between
archetypes.

Precision and Combat show moderate positive
associations:

\[
r_{\text{Precision}} = 0.376
\]

\[
r_{\text{Combat}} = 0.373
\]

These indicate that older observations within
these categories tend to have somewhat
higher BMI.

Endurance exhibits a negative association:

\[
r_{\text{Endurance}} = -0.308
\]

Team Sports shows almost no linear relationship:

\[
r_{\text{Team Sports}} = -0.043
\]

These differences suggest that there is no
single consistent age--BMI relationship
across all functional archetypes.

The observed correlations are descriptive
and cannot determine whether aging itself
causes changes in BMI.


\section{Limitations of the Analysis}

Several factors should be considered when
interpreting these results.

\begin{enumerate}[leftmargin=*]

\item \textbf{Unequal Sample Sizes:}

Archetype sizes range from 29 observations
in Technical / Mixed to 245 in Team Sports,
affecting the stability of descriptive estimates.

\item \textbf{Unequal Sex Representation:}

Male observations constitute 579 of 737
observations (78.6\%), which may influence
pooled anthropometric averages.

\item \textbf{Repeated Athletes:}

Some athletes appear in multiple Olympic
editions or archetypes.

The analysis retains one record per
athlete-year-archetype rather than one record
per unique athlete.

Therefore, observations are not necessarily
independent.

\item \textbf{Broad Archetype Classifications:}

Each archetype combines sports or events
with potentially different physical demands.

Aggregate results may conceal important
sport-specific variation.

\item \textbf{BMI Limitations:}

BMI does not distinguish lean muscle mass
from fat mass, making it an imperfect measure
of body composition among elite athletes.

\item \textbf{Dataset Scope:}

The conclusions apply to the recorded Indian
Olympic athlete observations with available
measurements and should not be generalized
to all Indian athletes.

\end{enumerate}


\section{Overall Interpretation and Conclusion}

The descriptive analysis reveals distinct
differences in age, height, weight, and BMI
across the seven functional archetypes
of Indian Olympic athletes.

Precision athletes exhibit the highest mean
age (29.21 years), weight (71.57 kg),
and BMI (24.07), while Racket Sports athletes
have the highest mean height (176.66 cm).

In contrast, Endurance athletes have the
lowest mean weight (61.04 kg) and BMI (20.36),
indicating a relatively lighter physical profile.

The analysis also reveals substantial differences
in variability.

Precision has the greatest age variability,
while Speed \& Power exhibits the greatest
variation in height and weight.

Combat and Precision show comparatively
large BMI dispersion, whereas Endurance
demonstrates a more concentrated BMI
distribution.

Sex-wise comparisons reveal that male
athletes have higher average height,
weight, and BMI across all archetypes.

However, some key patterns remain consistent
within both sexes.

Racket Sports athletes are the tallest
on average, and Endurance athletes exhibit
the lowest mean weight and BMI in both
male and female subgroups.

The height--weight analysis identifies a
positive association ($r = 0.698$),
while the age--BMI relationship is considerably
weaker ($r = 0.240$) and varies across
archetypes.

Overall, the findings suggest that the
functional categories differ in their
observed anthropometric profiles.

These patterns are consistent with the
possibility that different sports involve
different physical demands.

However, the descriptive analysis cannot
establish these demands as the cause of
the observed differences.

The findings should also be interpreted
in light of unequal sample sizes, sex
imbalance, repeated athlete observations,
and variation among sports within each
archetype.

\bigskip

\noindent
\textbf{Final Conclusion:}

Functional archetype classification provides
a useful descriptive framework for understanding
variation in the physical characteristics
of recorded Indian Olympic athletes.

The results highlight the importance of
examining both central tendencies and
within-group variability, as well as
accounting for sex composition and
sample-size differences when comparing
athletic populations.

\end{document}
